\documentclass{article}
\usepackage{amsmath,amssymb,amsthm}
\usepackage{algorithm}
\usepackage{algpseudocode}
\usepackage{hyperref}
\usepackage{enumitem}
\newtheorem{theorem}{Theorem}
\newtheorem{lemma}{Lemma}
\newtheorem{assumption}{Assumption}
\newcommand{\inner}[2]{\langle #1,#2\rangle}
\newcommand{\grad}{\nabla}
\newcommand{\prox}{\operatorname{prox}}
\newcommand{\Breg}[2]{D_h(#1,#2)}
\newcommand{\R}{\mathbb{R}}
\newcommand{\argmin}{\operatorname*{arg\,min}}

\begin{document}

\section*{Mirror Descent with Bregman Prox}

\section*{Mathematical Formulation}
Let $f:\R^d\rightarrow\R$ be a convex differentiable objective and let $h:\R^d\rightarrow\R$ be a \emph{prox‑function} that is $\sigma$–strongly convex with respect to a norm $\|\cdot\|$, i.e.
\[
h(y)\ge h(x)+\inner{\grad h(x)}{y-x}+\frac{\sigma}{2}\|y-x\|^{2},\qquad\forall x,y\in\R^d .
\]
The associated Bregman divergence is
\[
\Breg{x}{y}=h(x)-h(y)-\inner{\grad h(y)}{x-y}.
\]

Given a stepsize sequence $\{\eta_k\}_{k\ge0}$, Mirror Descent (MD) generates a sequence $\{x_k\}$ by
\begin{equation}\label{eq:MD}
x_{k+1}= \argmin_{x\in\R^d}\Big\{ \inner{\grad f(x_k)}{x}+ \frac{1}{\eta_k}\Breg{x}{x_k}\Big\}.
\end{equation}
When $h(x)=\tfrac12\|x\|_2^{2}$ the divergence reduces to $\tfrac12\|x-x_k\|_2^{2}$ and \eqref{eq:MD} collapses to the classical gradient step.

\section*{Pseudocode}
\begin{algorithm}[H]
\caption{Mirror Descent (MD)}\label{alg:MD}
\begin{algorithmic}[1]
\Require Convex $f$, prox‑function $h$, stepsizes $\{\eta_k\}$, initial point $x_0\in\operatorname{dom}h$.
\For{$k=0,1,\dots,T-1$}
    \State Compute gradient $g_k\gets\grad f(x_k)$.
    \State Solve
    \[
    x_{k+1}\gets \argmin_{x\in\R^d}\Big\{ \inner{g_k}{x}+ \frac{1}{\eta_k}\Breg{x}{x_k}\Big\}.
    \]
\EndFor
\State \textbf{return} $x_T$ (or the average $\bar x_T$).
\end{algorithmic}
\end{algorithm}

\section*{Dry Run Example}
Consider the one‑dimensional problem
\[
f(w)=(w-3)^{2},\qquad w_0=0,\qquad \eta_k\equiv\eta=0.1,
\]
and choose the Euclidean prox $h(w)=\tfrac12 w^{2}$, so that $\Breg{w}{w_k}= \tfrac12 (w-w_k)^{2}$.
Because $h$ is quadratic, the minimizer in \eqref{eq:MD} can be written explicitly:
\[
w_{k+1}= w_k-\eta\,\grad f(w_k).
\]

\begin{center}
\begin{tabular}{c|c|c|c}
$k$ & $w_k$ & $\grad f(w_k)=2(w_k-3)$ & $w_{k+1}=w_k-\eta\,\grad f(w_k)$ \\ \hline
0 & $0$ & $-6$ & $0-0.1(-6)=0.6$ \\
1 & $0.6$ & $2(0.6-3)=-4.8$ & $0.6-0.1(-4.8)=1.08$ \\
2 & $1.08$ & $2(1.08-3)=-3.84$ & $1.08-0.1(-3.84)=1.464$ \\
\end{tabular}
\end{center}
The corresponding objective values are $f(0)=9$, $f(0.6)=5.76$, $f(1.08)=3.6864$, $f(1.464)=2.3405$, illustrating monotone decrease.

\section*{Assumptions}
\begin{assumption}[Convexity and Smoothness]\label{as:convex}
$f$ is convex on $\R^d$ and $L$‑smooth w.r.t. the norm $\|\cdot\|$ induced by $h$, i.e.
\[
f(y)\le f(x)+\inner{\grad f(x)}{y-x}+\frac{L}{2}\|y-x\|^{2},\qquad\forall x,y\in\R^d .
\]
\end{assumption}
\begin{assumption}[Strong Convexity of $h$]\label{as:hstrong}
$h$ is $\sigma$‑strongly convex w.r.t. $\|\cdot\|$ (as defined above) with $\sigma>0$.
\end{assumption}
\begin{assumption}[Stepsize]\label{as:stepsize}
The stepsizes are constant $\eta_k\equiv\eta$ satisfying
\[
0<\eta\le \frac{\sigma}{L}.
\]
\end{assumption}

\section*{Main Convergence Theorem}
\begin{theorem}\label{thm:MD}
Let $\{x_k\}$ be generated by \eqref{eq:MD} under Assumptions~\ref{as:convex}--\ref{as:stepsize}. Denote $x^{\star}\in\arg\min f$ and let $D_{0}\triangleq \Breg{x^{\star}}{x_0}$. Then for any $T\ge1$,
\[
f(x_T)-f(x^{\star})\;\le\;\frac{D_{0}}{\eta\,T}.
\]
Consequently, MD attains an $O(1/T)$ convergence rate in function value.
\end{theorem}

\section*{Theoretical Properties}
\begin{itemize}[leftmargin=*]
\item \textbf{Stability.} The bound depends only on the initial Bregman distance $D_{0}$, not on the iterates themselves; thus MD is stable under any bounded $D_{0}$.
\item \textbf{Stepsize Sensitivity.} The admissible range $\eta\le\sigma/L$ guarantees that the term $\frac{1}{\eta}-\frac{L}{\sigma}$ appearing in the proof is non‑negative; larger $\eta$ may violate monotonicity.
\item \textbf{Comparison.} With the Euclidean prox ($\sigma=1$) the theorem recovers the classic gradient descent rate $O(L\|x_0-x^{\star}\|^{2}/T)$. For non‑Euclidean $h$ the bound scales with the Bregman distance, often yielding tighter results (e.g., KL‑divergence for simplex constraints).
\end{itemize}

\section*{Discussion}
The $O(1/T)$ rate matches the optimal deterministic rate for first‑order methods on convex smooth problems. By selecting a prox $h$ adapted to the geometry of the feasible set (e.g., entropy for probability simplices), MD can exploit structure that plain gradient descent cannot, leading to smaller $D_{0}$ and therefore faster practical convergence.

\section*{Preliminaries (Definitions and Lemmas)}
\begin{lemma}[Three‑Point Identity]\label{lem:3pt}
For any $x,y,z\in\R^d$,
\[
\Breg{x}{z}= \Breg{x}{y}+ \Breg{y}{z}+ \inner{\grad h(y)-\grad h(z)}{x-y}.
\]
\end{lemma}
\begin{proof}
Direct expansion of the definition of $D_h$ yields the identity. \hfill\qedsymbol
\end{proof}

\begin{lemma}[Optimality Condition of MD]\label{lem:optcond}
Let $x_{k+1}$ be defined by \eqref{eq:MD}. Then for every $x\in\R^d$,
\[
\inner{\grad f(x_k)}{x_{k+1}-x}+ \frac{1}{\eta_k}\inner{\grad h(x_{k+1})-\grad h(x_k)}{x_{k+1}-x}\le 0 .
\]
\end{lemma}
\begin{proof}
Since $x_{k+1}$ minimizes a convex function of $x$, the first‑order optimality condition gives
\[
0\in \grad f(x_k)+\frac{1}{\eta_k}\big(\grad h(x_{k+1})-\grad h(x_k)\big)+N_{\R^d}(x_{k+1}),
\]
which after taking inner product with any $x-x_{k+1}$ yields the claimed inequality. \hfill\qedsymbol
\end{proof}

\section*{Main Proof (Step‑by‑Step Derivation)}
We prove Theorem~\ref{thm:MD}. Throughout the proof we fix an arbitrary optimal point $x^{\star}\in\arg\min f$.

\bigskip
\noindent\textbf{[Step 1]} Apply Lemma~\ref{lem:optcond} with $x=x^{\star}$:
\[
\inner{\grad f(x_k)}{x_{k+1}-x^{\star}}+ \frac{1}{\eta}\inner{\grad h(x_{k+1})-\grad h(x_k)}{x_{k+1}-x^{\star}}\le 0 .
\tag{1}
\]

\noindent\textbf{[Step 2]} Rearrange (1) and add $\inner{\grad f(x_k)}{x_k-x_{k+1}}$ on both sides:
\[
\inner{\grad f(x_k)}{x_k-x^{\star}} \le
\inner{\grad f(x_k)}{x_k-x_{k+1}}-\frac{1}{\eta}\inner{\grad h(x_{k+1})-\grad h(x_k)}{x_{k+1}-x^{\star}} .
\tag{2}
\]

\noindent\textbf{[Step 3]} Use the three‑point identity Lemma~\ref{lem:3pt} with $(x,y,z)=(x^{\star},x_{k+1},x_k)$:
\[
\Breg{x^{\star}}{x_k}= \Breg{x^{\star}}{x_{k+1}}+ \Breg{x_{k+1}}{x_k}+ \inner{\grad h(x_{k+1})-\grad h(x_k)}{x^{\star}-x_{k+1}} .
\tag{3}
\]
Notice that the last inner product is the negative of the term appearing in (2).

\noindent\textbf{[Step 4] } Substitute the inner product from (3) into (2):
\[
\inner{\grad f(x_k)}{x_k-x^{\star}}
\le \inner{\grad f(x_k)}{x_k-x_{k+1}}-\frac{1}{\eta}\bigl[\Breg{x^{\star}}{x_k}-\Breg{x^{\star}}{x_{k+1}}-\Breg{x_{k+1}}{x_k}\bigr].
\tag{4}
\]

\noindent\textbf{[Step 5] } Expand $\inner{\grad f(x_k)}{x_k-x_{k+1}}$ using convexity of $f$:
\[
f(x_k)-f(x_{k+1})\le \inner{\grad f(x_k)}{x_k-x_{k+1}} .
\tag{5}
\]

\noindent\textbf{[Step 6] } Combine (4) and (5) and move $f(x_{k+1})$ to the left:
\[
f(x_k)-f(x^{\star}) \le
\frac{1}{\eta}\bigl[\Breg{x^{\star}}{x_k}-\Breg{x^{\star}}{x_{k+1}}-\Breg{x_{k+1}}{x_k}\bigr]
+ \bigl[f(x_k)-f(x_{k+1})\bigr] .
\tag{6}
\]

\noindent\textbf{[Step 7] } Cancel $f(x_k)$ on both sides of (6) to obtain:
\[
f(x_{k+1})-f(x^{\star}) \le
\frac{1}{\eta}\bigl[\Breg{x^{\star}}{x_k}-\Breg{x^{\star}}{x_{k+1}}-\Breg{x_{k+1}}{x_k}\bigr].
\tag{7}
\]

\noindent\textbf{[Step 8] } Use the $L$‑smoothness of $f$ (Assumption~\ref{as:convex}) to bound the Bregman term $\Breg{x_{k+1}}{x_k}$ from below. By definition of Bregman divergence and strong convexity of $h$ (Assumption~\ref{as:hstrong}),
\[
\Breg{x_{k+1}}{x_k}\ge \frac{\sigma}{2}\|x_{k+1}-x_k\|^{2}.
\tag{8}
\]
Moreover, smoothness yields
\[
f(x_{k+1})\le f(x_k)+\inner{\grad f(x_k)}{x_{k+1}-x_k}+\frac{L}{2}\|x_{k+1}-x_k\|^{2}.
\tag{9}
\]
Rearranging (9) gives
\[
\inner{\grad f(x_k)}{x_k-x_{k+1}}\le f(x_k)-f(x_{k+1})+\frac{L}{2}\|x_{k+1}-x_k\|^{2}.
\tag{10}
\]

\noindent\textbf{[Step 9] } From the optimality condition (Lemma~\ref{lem:optcond}) with $x=x_k$ we have
\[
\inner{\grad f(x_k)}{x_{k+1}-x_k} = -\frac{1}{\eta}\inner{\grad h(x_{k+1})-\grad h(x_k)}{x_{k+1}-x_k}
\le -\frac{\sigma}{\eta}\|x_{k+1}-x_k\|^{2},
\tag{11}
\]
where the inequality follows from the $\sigma$‑strong convexity of $h$ (i.e., monotonicity of its gradient).

\noindent\textbf{[Step 10] } Combine (10) and (11) to obtain
\[
f(x_k)-f(x_{k+1}) \ge \Big(\frac{\sigma}{\eta}-\frac{L}{2}\Big)\|x_{k+1}-x_k\|^{2}.
\tag{12}
\]
Because of the stepsize restriction $\eta\le\sigma/L$, the coefficient on the right‑hand side is non‑negative.

\noindent\textbf{[Step 11] } Drop the non‑negative term $\bigl(\frac{\sigma}{\eta}-\frac{L}{2}\bigr)\|x_{k+1}-x_k\|^{2}$ from (12) to obtain a simpler inequality:
\[
f(x_k)-f(x_{k+1}) \ge 0 .
\tag{13}
\]
Thus $f(x_k)$ is non‑increasing, and we may safely ignore the $\Breg{x_{k+1}}{x_k}$ term in (7) (it is non‑negative by (8)):

\[
f(x_{k+1})-f(x^{\star}) \le \frac{1}{\eta}\bigl[\Breg{x^{\star}}{x_k}-\Breg{x^{\star}}{x_{k+1}}\bigr].
\tag{14}
\]

\noindent\textbf{[Step 12] } Sum inequality (14) over $k=0,\dots,T-1$ telescopically:
\[
\sum_{k=0}^{T-1}\bigl[f(x_{k+1})-f(x^{\star})\bigr]
\le \frac{1}{\eta}\bigl[\Breg{x^{\star}}{x_0}-\Breg{x^{\star}}{x_T}\bigr]
\le \frac{D_{0}}{\eta},
\tag{15}
\]
where $D_{0}\triangleq\Breg{x^{\star}}{x_0}$ and we used $\Breg{x^{\star}}{x_T}\ge0$.

\noindent\textbf{[Step 13] } Since each term on the left of (15) is non‑negative, we obtain the bound for the last iterate:
\[
f(x_T)-f(x^{\star}) \le \frac{D_{0}}{\eta\,T}.
\tag{16}
\]
This completes the proof. \hfill\qedsymbol

\section*{Rate Derivation (With Constants)}
The constant in the $O(1/T)$ rate is explicitly $\frac{D_{0}}{\eta}$.
If the Euclidean prox is used ($h(x)=\frac12\|x\|_{2}^{2}$, $\sigma=1$) and $\eta=\frac{1}{L}$, then $D_{0}=\frac12\|x_0-x^{\star}\|_{2}^{2}$ and
\[
f(x_T)-f(x^{\star})\le \frac{L\,\|x_0-x^{\star}\|_{2}^{2}}{2T},
\]
the classic gradient‑descent bound.

\section*{Special Cases}
\begin{itemize}[leftmargin=*]
\item \textbf{Entropy prox on the simplex.}
Take $h(x)=\sum_{i}x_i\log x_i$ on the probability simplex. The induced Bregman distance is the Kullback–Leibler divergence. $D_{0}$ becomes $\mathrm{KL}(x^{\star}\|x_0)$, often much smaller than Euclidean distance, yielding faster practical convergence.
\item \textbf{Strongly convex $f$.}
If $f$ is $\mu$‑strongly convex, a similar derivation gives a linear rate $f(x_T)-f(x^{\star})\le (1-\eta\mu)^{T}D_{0}/\eta$.
\end{itemize}

\end{document}